\section{Discussion}

\subsection{Effect of CLAHE Preprocessing}

CLAHE did not improve Mean Test AUC in this study.
Experiment~02 (CLAHE + BCE) achieved AUC 0.8339 compared to the Experiment~01 baseline of 0.8352 ($-0.0013$).
Per-class effects were mixed: Edema improved by $+0.0065$ and Pleural Thickening by $+0.0092$, while Hernia decreased by $-0.0187$ and Fibrosis by $-0.0109$.
The inconsistency suggests that CLAHE enhances some local structures while potentially distorting others, and that its benefit is class-dependent.
Adding CBAM block34 on top of CLAHE (Exp~03b, AUC 0.8359) recovered some performance, but the CLAHE-free setup with the same architecture (Exp~07/08) ultimately performed better.
These results indicate that CLAHE preprocessing is not beneficial as a universal enhancement for this dataset and architecture.

\subsection{CBAM Placement}

The CBAM placement ablation (Experiments~03a, 03b, 03c) shows that block34 placement is most effective.
Placing CBAM only at block4 (03a) slightly reduced AUC relative to the no-CBAM baseline with CLAHE.
Using CBAM at all four blocks (block1234, 03c) achieved intermediate performance (0.8343), likely because inserting attention at early layers can over-attenuate low-level features before they contribute to high-level representations.
Block34, which adds attention at mid-to-high-level representations without intervening in early feature extraction, provides the best balance.

\subsection{Removing CLAHE and Using AdamW}

Experiment~07 (CBAM block34 + BCE + Adam, no CLAHE, AUC 0.8361) outperformed Experiment~03b (same but with CLAHE, AUC 0.8359), confirming that removing CLAHE is beneficial.
Switching from Adam to AdamW further improved Mean Test AUC from 0.8361 to 0.8371 (Exp~08, $+0.0010$).
AdamW's decoupled weight decay provides more consistent regularization during fine-tuning of pretrained features, which is the likely cause of this improvement.

\subsection{Effect of Loss Function: ASL vs. BCE}

The most notable finding is the behavior of Asymmetric Loss (Exp~09) relative to BCE (Exp~08).
The AUC difference is negligible ($-0.0001$), yet Recall improves by $+0.0990$ (0.2385 vs.\ 0.1395) and F1 improves by $+0.0774$ (0.2736 vs.\ 0.1962).

This behavior is consistent with how ASL operates: by applying a larger focusing weight to negative samples ($\gamma_- = 2.0$) and no down-weighting to positives ($\gamma_+ = 0.0$), ASL reduces the penalty for false negatives relative to BCE at the operating threshold.
The model trained with ASL is therefore more willing to predict positive at threshold 0.5, increasing Recall at the cost of Precision (0.3865 vs.\ 0.4482).

AUC and Recall/F1 measure different aspects of model behavior.
AUC evaluates the model's ability to rank positives above negatives across all thresholds, and is largely invariant to which threshold is chosen.
Recall at threshold 0.5 measures sensitivity at a specific operating point.
The fact that both models achieve nearly identical AUC (0.8370 vs.\ 0.8371) while differing substantially in Recall indicates that the ASL model redistributes confidence (shifts probabilities higher for positive cases) without fundamentally changing the AUC-level discrimination.
For clinical screening applications where missing a positive case is costly, Exp~09 is preferable; for ranking or prioritization tasks, both models are equivalent.

\subsection{Focal Loss}

Focal Loss with $\gamma{=}2.0$ (Exp~10, AUC 0.8351) did not improve over either BCE (0.8371) or ASL (0.8370).
Unlike ASL, Focal Loss applies the same focusing weight to both positive and negative samples, which may not provide the directional suppression of easy negatives needed in a heavily imbalanced multi-label setting.
Recall also decreased compared to BCE (0.1258 vs.\ 0.1395), suggesting that symmetric down-weighting with $\gamma{=}2.0$ did not selectively improve sensitivity.

\subsection{Grad-CAM Observations}

Grad-CAM visualizations were generated to qualitatively inspect model behavior.
The ASL model (Exp~09) consistently produces broader and more prominent activations on positive test examples, consistent with its higher predicted probabilities and recall-oriented behavior.
However, these observations are \emph{qualitative only}.
The visualizations should not be interpreted as evidence that the models correctly localize disease regions.
NIH ChestX-ray14 provides only image-level labels (no bounding boxes), so localization accuracy cannot be formally evaluated in this study.

\subsection{Limitations}

Several limitations should be noted.
First, NIH ChestX-ray14 labels were extracted via NLP and are known to contain noise, which may affect evaluation reliability.
Second, all threshold-dependent metrics are reported at a fixed threshold of 0.5; class-specific threshold tuning could improve these metrics.
Third, only frontal-view X-rays are considered.
Finally, model calibration was not evaluated; predicted probabilities may not be well-calibrated despite good AUC values.
